%!TEX root = ../template.tex
%%%%%%%%%%%%%%%%%%%%%%%%%%%%%%%%%%%%%%%%%%%%%%%%%%%%%%%%%%%%%%%%%%%
%% chapter1.tex
%% NOVA thesis document file
%%
%% Chapter with introduction
%%%%%%%%%%%%%%%%%%%%%%%%%%%%%%%%%%%%%%%%%%%%%%%%%%%%%%%%%%%%%%%%%%%

\typeout{NT FILE chapter1.tex}%

\chapter{Introduction}
\label{chapt:1}

\prependtographicspath{{Chapters/Figures/Covers/}}

% epigraph configuration
\epigraphfontsize{\small\itshape}
\setlength\epigraphwidth{12.5cm}
\setlength\epigraphrule{0pt}

\section{Basic Information}

\subsection{Definition and Importance of TSP}

The importance of the \Gls{TSP} goes beyond its seemingly simple formulation, branching into various fields and applications. In logistics and transportation,
solutions to the \Gls{TSP} enable optimal routing capable of leading to
significant cost savings and efficiency improvements. In
manufacturing, solution algorithms for \Gls{TSP} are used to optimize
drill movements in printed circuit board production or robotic
arms in assembly lines to reduce production times.
Beyond these, algorithms for \Gls{TSP} are employed in DNA
sequencing, astronomy, and even in the creation of works of art, showing
its versatility and importance in solving real-world problems.

\subsection{Historical Overview}

The Travelling Salesman Problem (\Gls{TSP}) represents one of the most studied challenges in computational mathematics and optimization, with a history that intertwines with the development of mathematical and operational research. Its evolution from early conceptual explorations to a crucial problem in algorithm design is a testament to its complexity and wide applicability.
The study of the \Gls{TSP} can be traced back to the works of Sir William Rowan Hamilton and Thomas Penyngton Kirkman in the 19th century, focused on Hamiltonian cycles and paths. Karl Menger, in the 1920s, introduced the "Messenger Problem", laying the foundation for the modern formulation of the \Gls{TSP} and emphasizing the search for efficient routing.

The advent of computers and operational research techniques marked a significant phase in \Gls{TSP} research. Early pioneers such as George Dantzig and Delbert Ray Fulkerson applied linear programming and cutting plane methods, illustrating the computational complexity of the problem and its potential for practical applications.

In the 1940s, the \Gls{TSP} began to find relevance in agriculture and statistics, with researchers using the problem's solution algorithms to optimize data collection and survey processes, showing the versatility of the \Gls{TSP} beyond theoretical mathematics.

In recent decades, the exploration of heuristic and bio-inspired algorithms, particularly the work of Marco Dorigo on combinatorial optimization through ant colonies, has opened new avenues for solving the \Gls{TSP}, highlighting continuous innovation in addressing its NP-hardness.


\section{The Problem}

\textbf{The Travelling Salesman Problem:} Given a set of cities and the distances between each pair of cities, find the shortest possible route that visits each city exactly once and returns to the original city.

\input{Chapters/Figures/USAmap.tex}

\begin{figure}[h]
	\centering
	\scalebox{0.5}{
		\begin{tikzpicture}
			% \tikzset{set state val/.style args={#1/#2}{#1={fill=blue!#2}}}
			\USA[every state={draw=white, ultra thick, fill=black!10}]

			\draw[green, ultra thick] (NY.center) -- (MI.center) -- (OH.center) -- (IL.center) -- (CO.center) -- (CA.center) -- (TX.center) -- (AL.center) -- (GA.center) -- (FL.center)  -- cycle;
			\draw[red, ultra thick] (NY.center) -- (OH.center) -- (GA.center) -- (MI.center) -- (IL.center) -- (TX.center) -- (CA.center)-- (CO.center) -- (AL.center) -- (FL.center)   -- cycle;


			% Randomly selected states
			\node[circle, fill=blue!50] at (CO) {CO};
			\node[circle, fill=blue!50] at (CA) {CA};
			\node[circle, fill=blue!50] at (NY) {NY};
			\node[circle, fill=blue!50] at (TX) {TX};
			\node[circle, fill=blue!50] at (AL) {AL};
			\node[circle, fill=blue!50] at (FL) {FL};
			\node[circle, fill=blue!50] at (GA) {GA};
			\node[circle, fill=blue!50] at (IL) {IL};
			\node[circle, fill=blue!50] at (MI) {MI};
			\node[circle, fill=blue!50] at (OH) {OH};

			% \draw[red, ultra thick] (CO.center) -- (CA.center) -- (NY.center) -- (TX.center) -- (AL.center) -- (FL.center) -- (GA.center) -- (IL.center) -- (MI.center) -- (OH.center) -- cycle;

			% Best TSP tour

			% Blue cycle
			% \draw[blue, ultra thick] (CO.center) -- (MI.center) -- (GA.center) -- (TX.center) -- (OH.center) -- (AL.center) -- (NY.center) -- (FL.center) -- (IL.center) -- (CA.center) -- cycle;
		\end{tikzpicture}
	}
	\caption{Two possible solutions for the (TSP) in a simplified instance of the United States, one shorter (in green) and one longer (in red).}
	\label{fig:us_map}
\end{figure}

\subsection{Mathematical Formulation}

Given a graph \(G = (V, E)\), where \(V\) is the set of vertices (cities) and \(E\) is the set of edges (paths between cities), with each edge \((i, j) \in E\) assigned a weight \(w_{ij}\) representing the cost of traveling from city \(i\) to city \(j\), the \Gls{TSP} seeks a Hamiltonian cycle of minimum weight.
The objective function to minimize is:
\[
	\min \sum_{i=1}^{n} \sum_{j \ne i, j=1}^{n} w_{ij} x_{ij},
\]
subject to constraints:
\[
	\sum_{i=1, i \ne j}^{n} x_{ij} = 1, \quad \forall j \in V,
\]
\[
	\sum_{j=1, j \ne i}^{n} x_{ij} = 1, \quad \forall i \in V,
\]
together with subtour elimination constraints, ensuring that each city is visited exactly once and the tour is cyclic.

\subsection{Integer Linear Programming Formulations}

\subsubsection{Miller–Tucker–Zemlin Formulation}

The \Gls{MTZ} formulation introduces auxiliary variables \(u_i\) for ordering, along with binary variables \(x_{ij}\):

Objective:
\[
	\min \sum_{i=1}^{n} \sum_{j \ne i, j=1}^{n} c_{ij} x_{ij},
\]
Subject to:
\[
	\sum_{i=1, i \ne j}^{n} x_{ij} = 1, \quad \forall j,
\]
\[
	\sum_{j=1, j \ne i}^{n} x_{ij} = 1, \quad \forall i,
\]
\[
	u_i - u_j + 1 \le (n-1)(1 - x_{ij}), \quad \forall 2 \le i \ne j \le n,
\]
\[
	2 \le u_i \le n, \quad \forall 2 \le i \le n.
\]

This formulation enforces a single tour through auxiliary ordering variables, effectively preventing subtours.

\subsubsection{Dantzig–Fulkerson–Johnson Formulation}

The DFJ formulation, known for its efficiency, directly introduces subtour elimination constraints:

Objective:
\[
	\min \sum_{i=1}^{n} \sum_{j \ne i, j=1}^{n} c_{ij} x_{ij},
\]
Subject to:
\[
	\sum_{i=1, i \ne j}^{n} x_{ij} = 1, \quad \forall j,
\]
\[
	\sum_{j=1, j \ne i}^{n} x_{ij} = 1, \quad \forall i,
\]
\[
	\sum_{i \in S, j \notin S} x_{ij} \ge 1, \quad \forall S \subset V, S \ne \emptyset, S \ne V.
\]


This formulation combats subtours by ensuring that at least one edge exits from every subset of vertices, guaranteeing a single connected tour.

% TODO: Fix this

% \begin{table}[h!]
% 	\centering
% 	\begin{tabularx}{\textwidth}{|X|X|X|}
% 			\hline
% 			\textbf{Caratteristica} & \textbf{Miller-Tucker-Zemlin (MTZ)} & \textbf{Dantzig-Fulkerson-Johnson (DFJ)} \\
% 			\hline
% 			\textbf{Funzione Obiettivo} & \(\min \sum_{i=1}^{n} \sum_{j \ne i, j=1}^{n} c_{ij} x_{ij}\) & \(\min \sum_{i=1}^{n} \sum_{j \ne i, j=1}^{n} c_{ij} x_{ij}\) \\
% 			\hline
% 			\textbf{Vincoli} & 
% 			\begin{tabular}{@{}l@{}}
% 			\(\sum_{i=1, i \ne j}^{n} x_{ij} = 1, \ \forall j\) \\
% 			\(\sum_{j=1, j \ne i}^{n} x_{ij} = 1, \ \forall i\) \\
% 			\(u_i - u_j + 1 \le (n-1)(1 - x_{ij}), \ \forall 2 \le i \ne j \le n\) \\
% 			\(2 \le u_i \le n, \ \forall 2 \le i \le n\)
% 			\end{tabular} & 
% 			\begin{tabular}{@{}l@{}}
% 			\(\sum_{i=1, i \ne j}^{n} x_{ij} = 1, \ \forall j\) \\
% 			\(\sum_{j=1, j \ne i}^{n} x_{ij} = 1, \ \forall i\) \\
% 			\(\sum_{i \in S, j \notin S} x_{ij} \ge 1, \ \forall S \subset V, S \ne \emptyset, S \ne V\)
% 			\end{tabular} \\
% 			\hline
% 			\textbf{Punti di Forza} & 
% 			\begin{tabular}{@{}l@{}}
% 			Riduce il problema a un singolo tour \\
% 			Previene efficacemente i sottoitinerari
% 			\end{tabular} & 
% 			\begin{tabular}{@{}l@{}}
% 			Include direttamente i vincoli di eliminazione dei sottoitinerari \\
% 			Efficiente per problemi di dimensioni piccole e medie
% 			\end{tabular} \\
% 			\hline
% 			\textbf{Limitazioni} & 
% 			\begin{tabular}{@{}l@{}}
% 			Numero elevato di variabili e vincoli \\
% 			Computazionalmente intensivo per problemi grandi
% 			\end{tabular} & 
% 			\begin{tabular}{@{}l@{}}
% 			Meno efficace per istanze molto grandi \\
% 			I vincoli possono essere ridondanti o costosi computazionalmente
% 			\end{tabular} \\
% 			\hline
% 	\end{tabularx}
% 	\caption{Confronto Trasposto delle Formulazioni del Problema del Viaggiatore (TSP)}
% 	\label{tab:tsp_formulations_transposed}
% \end{table}


\subsection{Variants and Applications}

The \Gls{TSP} adapts to a wide range of practical challenges, demonstrating its versatility. Key variants include:

\begin{itemize}
	\item \textbf{Symmetric vs. Asymmetric TSP (sTSP vs. aTSP):} sTSP implies that the weight of an edge between two nodes is equal regardless of the direction of travel, while aTSP allows distinct weights, reflecting the complexity of real-world routes.
	\item \textbf{Euclidean TSP:} This variant places cities in a Euclidean plane, emphasizing direct and straight-line distances, fundamental for drone navigation and infrastructure planning.
	\item \textbf{TSP with Time Windows:} Adds specific time intervals within which each city must be visited, crucial for planning in delivery services and maintenance operations.
	\item \textbf{Vehicle Routing Problem (VRP):} An extension of the \Gls{TSP} that involves multiple vehicles, central to optimizing fleet logistics and distribution networks.
\end{itemize}

The application of the \Gls{TSP} and its variants extends to various fields, showing its utility in addressing complex optimization problems:

\begin{itemize}
	\item \textbf{Logistics Optimization:} The Euclidean \Gls{TSP} and VRP are essential for reducing the cost of delivery routes and improving logistics efficiency.
	\item \textbf{Manufacturing:} The \Gls{TSP} helps optimize production processes, reducing downtime and improving overall workflows.
	\item \textbf{Planning and Scheduling:} The \Gls{TSP} with Time Windows ensures timely operations, optimizing resource allocation for services with critical timing requirements.
	\item \textbf{Network Design:} \Gls{TSP} algorithms facilitate the development of cost-effective and efficient telecommunication and transportation networks.
	\item \textbf{Genomics:} In genetics, the \Gls{TSP} supports advances in genome sequencing, offering insights into complex biological structures and processes.
\end{itemize}

% \begin{tikzpicture}
%   \path[mindmap,concept color=orange!40,fill=orange!20] 
%     node[concept] {TSP Variants and Applications}
%     [clockwise from=0]
%     child[grow=0,concept color=blue!50,fill=blue!20] 
%       node[concept] {Symmetric TSP}
%       child[grow=0,concept color=purple!50,fill=purple!20] 
%         node[concept] {Logistics}
%       child[grow=90,concept color=green!50,fill=green!20] 
%         node[concept] {Manufacturing}
%     child[grow=60,concept color=red!50,fill=red!20] 
%       node[concept] {Asymmetric TSP}
%       child[grow=0,concept color=cyan!50,fill=cyan!20] 
%         node[concept] {Telecommunications}
%       child[grow=90,concept color=yellow!50,fill=yellow!20] 
%         node[concept] {Network Design}
%     child[grow=120,concept color=orange!50,fill=orange!20] 
%       node[concept] {Euclidean TSP}
%       child[grow=0,concept color=gray!50,fill=gray!20] 
%         node[concept] {Drone Navigation}
%       child[grow=90,concept color=blue!50,fill=blue!20] 
%         node[concept] {Infrastructure Planning}
%     child[grow=180,concept color=green!50,fill=green!20] 
%       node[concept] {TSP with Time Windows}
%       child[grow=0,concept color=red!50,fill=red!20] 
%         node[concept] {Delivery Services}
%       child[grow=90,concept color=purple!50,fill=purple!20] 
%         node[concept] {Maintenance Scheduling};
% \end{tikzpicture}

\section{Thesis Objectives}



This thesis aims to explore the Travelling Salesman Problem (TSP) from both a theoretical and practical perspective. The focus is on examining the Red-Black Ant Colony System \Gls{RB-ACS}, a variant of the Ant Colony System algorithm \Gls{ACS}, with some modifications, to understand its effectiveness in addressing the challenges of the \Gls{TSP}. The study has two objectives: improving the understanding of the \Gls{TSP} in computational optimization and evaluating the performance of the modified \Gls{RB-ACS} compared to conventional solutions.

\subsection{Main Objectives}

The thesis is guided by the following objectives:

\begin{itemize}
	\item Examine the existing body of work on the \Gls{TSP}, gathering its mathematical foundations, evolution, and the set of strategies previously devised for its resolution.
  \item Detail the operation of the \Gls{RB-ACS}, emphasizing its theoretical foundations, operational mechanics, and its distinctions from more traditional approaches.
  \item Conduct a systematic experimental comparison of the \Gls{RB-ACS} against standard \Gls{TSP} solution algorithms across a range of problem sets, aiming to discern its operational effectiveness and scalability.
	\item Identify practical scenarios where the modified algorithm could be applied, showing its relevance and potential benefits in tangible contexts.
\end{itemize}

% TODO: this is nice, but needs fixing
% \begin{figure}[h!]
% 	\centering
% 	\begin{tikzpicture}[node distance=2cm, every node/.style={font=\small}]

% 		% Define the styles for the nodes
% 		\tikzstyle{startstop} = [rectangle, rounded corners, minimum width=3cm, minimum height=1cm,text centered, draw=black, fill=red!30]
% 		\tikzstyle{process} = [rectangle, minimum width=3cm, minimum height=1cm, text centered, draw=black, fill=orange!30]
% 		\tikzstyle{decision} = [diamond, minimum width=3cm, minimum height=1cm, text centered, draw=black, fill=green!30]
% 		\tikzstyle{arrow} = [thick,->,>=stealth]

% 		% Nodes
% 		\node (start) [startstop] {Start};
% 		\node (exact) [process, below of=start] {Exact Methods};
% 		\node (heuristic) [process, right of=exact, xshift=4cm] {Heuristic Methods};
% 		\node (branchbound) [process, below of=exact] {Branch and Bound};
% 		\node (ip) [process, below of=branchbound] {Integer Programming};
% 		\node (sa) [process, below of=heuristic] {Simulated Annealing};
% 		\node (ga) [process, below of=sa] {Genetic Algorithms};
% 		\node (end) [startstop, below of=ip] {End};

% 		% Arrows
% 		\draw [arrow] (start) -- (exact);
% 		\draw [arrow] (start) -- (heuristic);
% 		\draw [arrow] (exact) -- (branchbound);
% 		\draw [arrow] (exact) -- (ip);
% 		\draw [arrow] (heuristic) -- (sa);
% 		\draw [arrow] (heuristic) -- (ga);
% 		\draw [arrow] (branchbound) -- (end);
% 		\draw [arrow] (ip) -- (end);
% 		\draw [arrow] (sa) -- (end);
% 		\draw [arrow] (ga) -- (end);

% 	\end{tikzpicture}
% 	\caption{Flowchart of Exact and Heuristic Methods for Solving TSP}
% 	\label{fig:tsp_flowchart}
% \end{figure}
%\subsection{Ambito dell'Indagine}
%
%La ricerca approfondisce aree di interesse specifiche:
%
%\begin{itemize}
	%\item Una panoramica analitica del \Gls{TSP}, che copre le sue definizioni, implicazioni e varianti, gettando le basi per le successive indagini algoritmiche.
	%\item Una revisione approfondita degli algoritmi euristici e meteuristici, concentrandosi in particolare sull'Ottimizzazione della Colonia di Formiche come fondamento per il Sistema di Colonie di Formiche Red-Black.
%	\item L'adattamento e l'applicazione del Sistema di Colonie di Formiche Red-Black, illustrando gli aspetti innovativi introdotti e le loro giustificazioni teoriche.
%	\item Un confronto complessivo delle prestazioni degli algoritmi, utilizzando diverse istanze di \Gls{TSP} per valutare l'efficacia e la qualità delle soluzioni.
%\end{itemize}

%Attraverso questa indagine, la tesi si sforza di offrire contributi significativi allo studio del \Gls{TSP}, dimostrando l'utilità del Sistema di Colonie di Formiche Red-Black nella risoluzione di problemi di ottimizzazione e aprendo la strada a future ricerche.

%\section{Ambito e Limitazioni}

% Questa tesi di laurea si muove attraverso specifiche definizioni di ambito e incontra limitazioni cruciali da delineare per un'interpretazione accurata dei risultati dello studio e della loro rilevanza per il campo più ampio.

\subsection{Theoretical Constraints}

The exploration within this thesis is tailored to fit the available computational resources, focusing on \Gls{TSP} instances that balance presenting a notable challenge and the feasibility of analysis within the constraints of our computational setup. A notable aspect of this research concerns the reimplementation in Rust of all algorithms discussed, due to the unavailability of the original source code for some of these algorithms. This necessity highlights a possible variation in performance metrics, which might not fully align with those derived from the authors' original implementations. Such differences could stem from various factors, including the intrinsic characteristics of Rust as a programming language and the complexities involved in translating complex algorithmic logic into code. Consequently, while aiming to provide a comparative analysis of algorithmic performance, it is imperative to consider these reimplementation factors as potential variables influencing the results.

Furthermore, while acknowledging the potential for real-world applications of the algorithms under study, in-depth case studies that go beyond theoretical exploration and computational experimentation fall outside the scope of this bachelor's thesis. The main focus remains on the theoretical foundations and computational evaluation of \Gls{TSP} solutions.

\subsection{Practical Considerations}

The practicality of this thesis is influenced by several limitations that guide its experimental approach and evaluation methodology:

\begin{itemize}
  \item The performance evaluation of the \Gls{RB-ACS} is conducted through a carefully chosen subset of \Gls{TSP} instances. These instances aim to provide a diverse but not exhaustive representation of potential problem scenarios, highlighting the intrinsic limitation in capturing the entire spectrum of \Gls{TSP} configurations.
	\item Due to computational resource constraints, the scalability of the proposed algorithm is tested in a constrained environment. This limitation may hinder the extrapolation of performance results to significantly larger or more complex \Gls{TSP} instances.
	\item The thesis considers the practical implications of its findings in a theoretical and computational context, without delving into field testing or specific case studies for applications.
\end{itemize}

%Queste considerazioni sono fondamentali per modellare l'approccio della ricerca e interpretare i risultati di questa tesi di laurea, stabilendo aspettative realistiche per i suoi contributi alla comprensione e alla risoluzione del Problema del Viaggiatore.

\section{Thesis Structure}

\subsection{Chapter Overview}
Chapter \ref{chapt:1} introduces the Travelling Salesman Problem (TSP), covering what it is, why it is important, and a bit of its history. It outlines the main questions this thesis aims to answer and briefly outlines the methods used.
Chapter \ref{chapt:2} analyzes the overall framework of computational complexity, focusing on the famous P vs. NP question, to help understand why solving the \Gls{TSP} can be so challenging.
The discussion continues in Chapters \ref{chapt:3} and \ref{chapt:4} to explore different ways of approaching the \Gls{TSP}. Chapter \ref{chapt:3} examines exact solutions while Chapter \ref{chapt:4} discusses heuristic and metaheuristic methods, including well-known strategies such as genetic algorithms and simulated annealing.
Chapter \ref{chapt:5} focuses on the \Gls{ACS}, explaining how it works, why it is nature-inspired, and how it applies to solving the \Gls{TSP}.
Then, Chapter \ref{chapt:6} introduces a variant of the \Gls{ACS}, the \Gls{RB-ACS}. This chapter explains how it was developed, the new ideas it brings, and how tests demonstrate that it could offer an advantage in solving \Gls{TSP} instances.
% Il Capitolo~\ref{chapt:7} presenta "Ibn-Battuta", uno strumento creato per visualizzare come diversi algoritmi \Gls{TSP}, incluso il Sistema di Colonie di Formiche Red-Black, risolvono i problemi. Questo rende più facile vedere come funzionano gli algoritmi nella pratica
Chapter ~\ref{chapt:7} presents the experimental results obtained.
Finally, Chapter \ref{chapt:8} concludes the thesis by summarizing the highlights of this thesis.
This thesis aims to contribute to the conversation on solving the \Gls{TSP}, seeking to bridge the gap between theoretical challenges and real-world solutions, and opening the door to further investigations into optimization problems.
